% ---------------------------------------------------
% Trabajo Final de Grado
% Author: Fabián González Lence <alu0101549491@ull.edu.es>
% Chapter: Desarrollo
% ----------------------------------------------------

\chapter{Desarrollo} \label{chap:Development}

Cada uno de los cinco proyectos recorrió el flujo en cuatro fases descrito en la Sección~\ref{sec:flujo-trabajo} partiendo de los artefactos de diseño elaborados en el Capítulo~\ref{chap:Design}.
Este capítulo documenta la ejecución concreta de ese flujo: qué modelo ejerció de responsable en cada fase, qué decisiones se tomaron durante el proceso, qué incidencias surgieron y cómo se resolvieron antes de avanzar a la fase siguiente.

\section{The Hangman Game} \label{sec:dev-hangman}

\ac{HANGMAN} fue el primer proyecto en completar el flujo de cuatro fases y actuó como banco de calibración de las plantillas de \textit{prompts}: su alcance acotado —once ficheros de producción, sin \textit{framework}— permitió aislar el comportamiento de cada modelo antes de escalar en complejidad.

\subsection{Fase de Arquitectura}

El \textit{prompt} de arquitectura\footnote{\href{https://github.com/alu0101549491/TFG-Fabian-Gonzalez-Lence/blob/main/resources/prompts/1-TheHangmanGame/1-Inicializacion/architecture.prompt.md}{\textit{Prompt} de arquitectura}} proporcionó a Claude los tres artefactos de diseño del Capítulo~\ref{chap:Design}, junto con la pila tecnológica acordada y un conjunto de reglas de estilo de código, y solicitó la generación del árbol de ficheros con esqueletos de clases documentados y la configuración inicial del proyecto. Claude devolvió en una única interacción la estructura completa del proyecto\footnote{\href{https://github.com/alu0101549491/TFG-Fabian-Gonzalez-Lence/blob/main/resources/project-structures/1-TheHangmanGame.md}{Estructura de código de \texttt{HANGMAN}}}: la estructura de directorios reproducía con exactitud las capas \ac{MVC} y la jerarquía de componentes de vista definidas en los diagramas, y la compilación resultó limpia en el primer intento. La estructura de ficheros se trasladó al repositorio manualmente, ya que el modelo la proporcionó como texto. Este resultado validó la elección de Claude como Arquitecto y fijó el estándar de calidad esperado en los proyectos siguientes.\\

\subsection{Fase de Codificación}

Mistral implementó cada módulo de forma independiente tomando como contexto el \textit{prompt} base de codificación\footnote{\href{https://github.com/alu0101549491/TFG-Fabian-Gonzalez-Lence/tree/main/resources/prompts/1-TheHangmanGame/3-CodigoPorModulo}{\textit{Prompts} de código}}, que aportaba la especificación de requisitos, el diagrama de clases y el esqueleto de la clase correspondiente. El \texttt{index.html} y la hoja de estilos fueron los únicos artefactos integrados directamente sin modificación, señalando una fortaleza del modelo en código de presentación estático. La incidencia más destacable surgió en \textit{WordDictionary}: en lugar de externalizar el vocabulario, Mistral lo embebió como un \textit{array} literal dentro del método \textit{initializeAnimalWords()}\footnote{\href{https://github.com/alu0101549491/TFG-Fabian-Gonzalez-Lence/blob/6ade2d937c5fc8507893551a8f82d8d46e6fcbb2/projects/1-TheHangmanGame/src/models/word-dictionary.ts\#L46-L54}{\texttt{word-dictionary.ts}, líneas 46--54} — versión inicial con \textit{array} embebido}. La solución era funcional, pero introducía deuda de mantenibilidad —ampliar el diccionario exigía modificar el fuente— y se registró como incidencia abierta para la fase siguiente.

\subsection{Fase de Revisión}

GitHub Copilot revisó los componentes de forma secuencial aplicando la plantilla de revisión\footnote{\href{https://github.com/alu0101549491/TFG-Fabian-Gonzalez-Lence/tree/main/resources/prompts/1-TheHangmanGame/4-RevisionPorModulo}{\textit{Prompts} de revisión}}, que aporta como contexto la especificación, el diagrama de clases y el código a evaluar, y produce un veredicto estructurado por requisito. Ante la incidencia de \textit{WordDictionary}, el Revisor no se limitó a señalarla: generó autónomamente \texttt{src/data/animals.json} y reescribió el constructor para cargarlo mediante importación estática\footnote{\href{https://github.com/alu0101549491/TFG-Fabian-Gonzalez-Lence/blob/main/projects/1-TheHangmanGame/src/models/word-dictionary.ts\#L30-L34}{\texttt{word-dictionary.ts}, líneas 30--34} — versión final con vocabulario externalizado}, cerrando la incidencia. Este comportamiento, fue el primer indicio de que el Revisor podía absorber correcciones de implementación, capacidad que el otro enfoque explotaría a partir de \ac{CARTO} (Sección~\ref{sec:approach-change}). El resto de módulos recibieron un veredicto aprobado en su primera revisión.

\subsection{Fase de Pruebas}

Qwen AI generó nueve \textit{suites} de pruebas unitarias a partir de aplicar la plantilla de \textit{testing}\footnote{\href{https://github.com/alu0101549491/TFG-Fabian-Gonzalez-Lence/tree/main/resources/prompts/1-TheHangmanGame/5-TestingPorModulo}{\textit{Prompts} de \textit{testing}}}, distribuidas en las tres capas \ac{MVC}; siete compilaron y pasaron sin correcciones. Las dos excepciones compartían raíz común: dependencia del entorno de Canvas bajo JSDOM. \texttt{hangman-renderer.test.ts} carecía del \textit{mock} del contexto 2D\footnote{\href{https://github.com/alu0101549491/TFG-Fabian-Gonzalez-Lence/blob/main/projects/1-TheHangmanGame/tests/views/hangman-renderer.test.ts\#L28-L45}{\texttt{hangman-renderer.test.ts}, líneas 28--45} — \textit{mock} del contexto 2D de Canvas}, y \texttt{game-view.test.ts} llamaba a \textit{initialize()} antes de limpiar los contadores de \textit{mock}, contaminando las aserciones\footnote{\href{https://github.com/alu0101549491/TFG-Fabian-Gonzalez-Lence/blob/main/projects/1-TheHangmanGame/tests/views/game-view.test.ts\#L56-L67}{\texttt{game-view.test.ts}, líneas 56--67} — orden de \textit{clearAllMocks()} e \textit{initialize()} en \textit{beforeEach}}. Ambas se resolvieron en una segunda iteración, y la cobertura final alcanzó el 82\% de líneas en todo el código.\\

% DEJAR ESTO PARA EL CAPITULO 8
\begin{comment}
El informe de calidad completo, elaborado sobre los criterios ISO/IEC 25010, puede consultarse en \texttt{projects/1-TheHangmanGame/docs/HANGMAN\_quality\_report.md}; su análisis cuantitativo se recoge en el Capítulo~\ref{chap:Results}.\\
\end{comment}